% Part: second-order-logic
% Chapter: syntax-and-semantics
% Section: size-of-domain

\documentclass[../../../include/open-logic-section]{subfiles}

\begin{document}

\olfileid{sol}{syn}{siz}

\olsection{अनंत आणि \usetoken{S}{enumerable}
  \usetoken{P}{domain} यांचे वर्णन}

%READERNOTE{SOL6-A092: Inf डेडेकिंड अनंतता व्यक्त करते. येथे सर्वसाधारण अनंततेशी आणि Fin च्या सांततेशी सममूल्यत्वासाठी निवडीचे स्वयंसिद्धक मानणारा सामान्य संचसैद्धान्तिक संदर्भ घेतला आहे. Orbit ची पुनरावृत्ती काढल्यावर repetition-free प्रगणना मिळते.}
येथे निवडीचे स्वयंसिद्धक मानणारा नेहमीचा
संचसैद्धान्तिक संदर्भ घेतला आहे. त्यात डेडेकिंड अनंतता
आणि सर्वसाधारण अनंतता सममूल्य असतात.
संच~$M$ हा
(डेडेकिंडच्या अर्थाने)
अनंत असतो,
जर आणि फक्त जर
$f\colon M
\to M$ असे
!!a{injective}
फलन असते जे
!!{surjective}
नसते; म्हणजे
$\ran{f} \neq M$.
प्रथम-क्रम
तर्कशास्त्रात
एक-स्थानी
!!{function}~$f$
घेऊन
!!a{structure}~$\Struct{M}$
मध्ये त्याला
नेमलेले फलन
$\Assign{f}{M}$
हे !!{injective}
आहे आणि
$\ran{f} \neq
\Domain{M}$
आहे असे
म्हणता येते:
\[
\lforall[x][\lforall[y][(\eq[f(x)][f(y)] \to \eq[x][y])]] \land
\lexists[y][\lforall[x][\eq/[y][f(x)]]].
\]
$\Struct{M}$
ही या
!!{sentence}
ची पूर्ति करत
असेल, तर
$\Assign{f}{M}:
\Domain{M} \to \Domain{M}$
हे !!{injective}
असते, म्हणून
$\Domain{M}$
अनंत असलेच
पाहिजे.
$\Domain{M}$
अनंत असेल
आणि म्हणून
असे फलन
अस्तित्वात असेल,
तर ते फलन
$\Assign{f}{M}$
म्हणून घेता
येते आणि
$\Struct{M}$
ही त्या
!!{sentence}
ची पूर्ति करेल.
मात्र यासाठी
आपल्या भाषेत
या हेतूने
वापरण्यात येणारे
अतार्किक चिन्ह~$f$
असणे आवश्यक
आहे.
द्वितीय-क्रम
तर्कशास्त्रात
असे फलन
\emph{अस्तित्वात आहे}
असे थेट
म्हणता येते.
यासाठी
आता~$f$
लागत नाही,
आणि शुद्ध
द्वितीय-क्रम
तर्कशास्त्रातील
पुढील
!!{sentence}
मिळते:
\[
\fn{Inf} \ident \lexists[u][(\lforall[x][\lforall[y][(\eq[u(x)][u(y)]
      \lif \eq[x][y])]] \land \lexists[y][\lforall[x][\eq/[y][u(x)]]])].
\]
$\Sat{M}{\fn{Inf}}$
हे
$\Domain{M}$
अनंत असेल
तर आणि तरच
खरे असते.
मग
$\fn{Fin} \ident \lnot \fn{Inf}$
असे परिभाषित
करता येते;
$\Sat{M}{\fn{Fin}}$
हे
$\Domain{M}$
सांत असेल
तर आणि तरच
खरे असते.
शुद्ध प्रथम-क्रम
तर्कशास्त्रातील
कोणतेही एक
!!{sentence}
!!{domain}
अनंत आहे
हे व्यक्त
करू शकत
नाही, मात्र
अशा वाक्यांचा
अनंत संच
ते करू शकतो.
शुद्ध प्रथम-क्रम
तर्कशास्त्रातील
!!{sentence}s
चा असा
कोणताही संच
नाही की
त्याची
!!a{structure}
मध्ये पूर्ति
होते जर आणि
फक्त जर तिचे
आधारक्षेत्र
सांत आहे.

\begin{prop}
$\Sat{M}{\fn{Inf}}$
हे
$\Domain{M}$
अनंत असेल
तर आणि तरच
खरे असते.
\end{prop}

\begin{proof}
$\Sat{M}{\fn{Inf}}$
असेल तर आणि
तरच काही~$s$
साठी
  $\Sat{M}{\lforall[x][\lforall[y][(\eq[u(x)][u(y)] \lif \eq[x][y])]]
    \land \lexists[y][\lforall[x][\eq/[y][u(x)]]]}[s]$.
  असे असेल,
  तर $s(u)$
  हे !!a{injective}
  फलन आहे
  आणि काही
  $y
  \in \Domain{M}$
  हे~$s(u)$
  च्या व्याप्तीत
  नाही.
  उलट
  $f\colon \Domain{M} \to \Domain{M}$
  असे
  !!a{injective}
  फलन असेल
  की
  $\ran{f}
  \neq \Domain{M}$,
  तर
  $s(u) = f$
  असे चर-मूल्यांकन
  मिळते.
\end{proof}

संच $M$
!!{enumerable}
असतो, जर
त्यातील
!!{element}s
ची पुनरावृत्ती
नसलेली
(पण शक्यतो
सांत असलेली)
पुढील
प्रगणना असते:
\[
m_0, m_1, m_2, \dots
\]
अशी प्रगणना
असते जर
आणि फक्त
जर
!!a{element}
$z \in M$
आणि
$f\colon M \to M$
असे फलन
असते की
$z$, $f(z)$,
$f(f(z))$, \dots,
हे~$M$
चे सर्व
!!{element}s
असतात.
कारण प्रगणना
असेल, तर
$z = m_0$
आणि
$f(m_k) = m_{k+1}$
(किंवा
$m_k$ हा
प्रगणनेतील
शेवटचा
!!{element}
असेल, तर
$f(m_k) = m_k$)
हे अपेक्षित
!!{element}
आणि फलन
ठरतात.
उलट असे
$z$ आणि
$f$ असतील,
तर
$z$, $f(z)$,
$f(f(z))$, \dots,
या यादीतील पुनरावृत्ती
काढल्यावर~$M$ ची प्रगणना
मिळते आणि
$M$ हा
!!{enumerable}
असतो.
$z$ आणि~$f$
यांचे अस्तित्व
द्वितीय-क्रम
तर्कशास्त्रात
व्यक्त करून
असे
!!{sentence}
तयार करता
येते की
ते
!!a{structure}
मध्ये खरे
असते जर
आणि फक्त
जर ती
!!{structure}
!!{enumerable}
असते:
\[
\fn{Count} \ident
\lexists[z][\lexists[u][\lforall[X][((X(z) \land
      \lforall[x][(X(x) \lif X(u(x)))]) \lif \lforall[x][X(x)])]]]
\]

\begin{prop}
$\Sat{M}{\fn{Count}}$
हे
$\Domain{M}$
!!{enumerable}
असेल तर
आणि तरच
खरे असते.
\end{prop}

\begin{proof}
$\Domain{M}$
!!{enumerable}
आहे असे
समजू आणि
$m_0$, $m_1$,
\dots, ही
प्रगणना घेऊ.
पुनरावृत्ती
काढून टाकल्यास
कोणतेही
$m_k$ दोनदा
येणार नाही
याची खात्री
करता येते.
$f(m_k) = m_{k+1}$
असे परिभाषित
करू (सांत
प्रगणनेच्या
शेवटच्या घटकावर
फलन तोच घटक
परत देते)
आणि
$s(z) = m_0$,
$s(u) = f$
घेऊ.
आता पुढील
गोष्ट दाखवू:
\[
\Sat{M}{\lforall[X][((X(z) \land \lforall[x][(X(x) \lif X(u(x)))])
    \lif \lforall[x][X(x)])]}[s]
\]
$M \subseteq \Domain{M}$
हा कोणताही
संच असू
द्या.
तसेच
$\Sat{M}{(X(z) \land \lforall[x][(X(x) \lif
X(u(x)))])}[\Subst{s}{M}{X}]$
असे धरू.
मग
$\Subst{s}{M}{X}(z) \in M$
आणि
$x \in M$
असेल तेव्हा
$(\Subst{s}{M}{X}(u))(x) \in M$
ही.
दुसऱ्या
शब्दांत,
$\varAssign{\Subst{s}{M}{X}}{s}{X}$
असल्यामुळे
$m_0 \in M$
आणि
$x \in M$
असेल तर
$f(x) \in M$;
म्हणून
$m_0 \in M$,
$m_1 = f(m_0) \in M$,
$m_2 = f(f(m_0)) \in M$
इत्यादी.
त्यामुळे
$M = \Domain{M}$
आणि म्हणून
$\Sat{M}{\lforall[x][X(x)]}[\Subst{s}{M}{X}]$.
$M \subseteq
\Domain{M}$
कोणताही
असल्यामुळे
सिद्धता पूर्ण:
$\Sat{M}{\fn{Count}}$.

आता
$\Sat{M}{\fn{Count}}$,
म्हणजे
\[
\Sat{M}{\lforall[X][((X(z) \land \lforall[x][(X(x) \lif X(u(x)))])
    \lif \lforall[x][X(x)])]}[s]
\]
काही~$s$
साठी आहे
असे समजू.
$m = s(z)$
आणि
$f = s(u)$
घेऊ आणि
$M = \{m,
f(m), f(f(m)), \dots\}$
हा संच
विचारात
घेऊ.
असा
परिभाषित
केलेला
$M$
स्पष्टपणे
!!{enumerable}
आहे.
मग
\[
\Sat{M}{(X(z) \land \lforall[x][(X(x) \lif X(u(x)))])
    \lif \lforall[x][X(x)]}[\Subst{s}{M}{X}]
\]
हे गृहीतकावरून
मिळते.
तसेच
$\Sat{M}{X(z)}[\Subst{s}{M}{X}]$,
कारण
$M \ni m =
\Subst{s}{M}{X}(z)$;
आणि
$\Sat{M}{\lforall[x][(X(x) \lif
X(u(x)))]}[\Subst{s}{M}{X}]$
सुद्धा,
कारण
$x \in M$
असेल तेव्हा
$f(x) \in
M$.
म्हणून
पूर्वपक्ष
आणि शर्त
दोन्हींची
पूर्ति
झाल्याने
उत्तरपक्षाचीही
पूर्ति झाली
पाहिजे:
$\Sat{M}{\lforall[x][X(x)]}[\Subst{s}{M}{X}]$.
पण याचा
अर्थ
$M = \Domain{M}$.
म्हणून
$\Domain{M}$
!!{enumerable}
आहे, कारण
$M$ व्याख्येने
तसाच आहे.
\end{proof}

\begin{prob}
!!{sentence}
$\fn{Inf} \land \fn{Count}$
हे नेमक्या
!!{denumerable}
आधारक्षेत्रांत
खरे असते.
$\fn{Count}$
ची व्याख्या
बदलून
आधारक्षेत्र
!!{denumerable}
आहे हे
थेट व्यक्त
करणारे
वेगळे
!!{sentence}
मिळवा आणि
ते तसे
करते हे
सिद्ध करा.
\end{prob}

\end{document}
